\chapter{反常积分}
\label{chap:improper-integrals-solutions}

\begin{solution}{improper-p-integrals-complete}
若 \(p\ne1\)，
\[
 \int_1^R x^{-p}\,dx=\frac{R^{1-p}-1}{1-p}.
\]
当 \(p>1\) 时趋于 \(1/(p-1)\)，当 \(p<1\) 时趋于 \(+\infty\)。
\(p=1\) 时积分为 \(\log R\)，发散。

类似地，若 \(p\ne1\)，
\[
 \int_\varepsilon^1x^{-p}\,dx
 =\frac{1-\varepsilon^{1-p}}{1-p}.
\]
当 \(p<1\) 时趋于 \(1/(1-p)\)，当 \(p>1\) 时发散；\(p=1\) 时为
\(-\log\varepsilon\)，也发散。
\end{solution}

\begin{solution}{improper-split-independence}
设在分割点 \(c\) 两侧都收敛，另取 \(d>c\)。则
\[
 \int_{-\infty}^{d}f
 =\int_{-\infty}^{c}f+\int_c^d f,
\qquad
 \int_d^\infty f
 =\int_c^\infty f-\int_c^d f.
\]
两式相加时有限区间积分消去，故总值不变。\(d<c\) 同理。
\end{solution}

\begin{solution}{improper-principal-value-not-integral}
函数为奇函数，故
\[
 \int_{-R}^{R}\frac{x}{1+x^2}\,dx=0,
\]
对称主值为零。但
\[
 \int_0^R\frac{x}{1+x^2}\,dx
 =\frac12\log(1+R^2)\to+\infty,
\]
左半轴积分趋于 \(-\infty\)。两侧都没有有限极限，所以通常反常积分发散。
\end{solution}

\begin{solution}{improper-interior-singularity-split}
\[
 \int_{-1}^{-\varepsilon}\frac{dx}{x}=\log\varepsilon\to-\infty,
\qquad
 \int_{\varepsilon}^{1}\frac{dx}{x}=-\log\varepsilon\to+\infty.
\]
两侧均发散，故反常积分不存在；但对称截断之和恒为零，所以主值为零。
\end{solution}

\begin{solution}{improper-cauchy-tail-sequence}
若 Cauchy 准则成立，任意 \(u_n,v_n\to\infty\) 且 \(u_n<v_n\) 最终
进入相应尾部，故积分趋零。反之，若 Cauchy 准则失败，则存在
\(\varepsilon_0>0\)，对每个整数 \(n\) 可选
\(v_n>u_n\ge n\) 使
\[
 \abs{\int_{u_n}^{v_n}f}\ge\varepsilon_0.
\]
这给出违背序列条件的两列。
\end{solution}

\begin{solution}{improper-log-power-infinity}
令 \(t=\log x\)，则 \(dx=e^t dt\)，积分化为
\[
 \int_1^\infty e^{-(p-1)t}t^{-q}\,dt.
\]
若 \(p>1\)，指数衰减可控制幂因子的增长，积分收敛；若 \(p<1\)，
被积函数最终大于1，故积分发散。若 \(p=1\)，得到
\(\int_1^\infty t^{-q}dt\)，恰在 \(q>1\) 时收敛。

具体地，对 \(c>0\)，取整数 \(m>\abs q\)。有限 Taylor 公式及
正余项给出 \(e^{ct/2}\ge(ct/2)^m/m!\)。因 \(m>\abs q\)，
充分大时右端大于 \(t^{\abs q}\)。当 \(c=p-1>0\) 时，遂有
\[
 e^{-ct}t^{-q}\le e^{-ct/2}.
\]
右端可积。若 \(p<1\)，令 \(c=1-p>0\)，则
同一估计给出 \(e^{ct}t^{-q}\ge e^{ct/2}\to+\infty\)，所以尾部
可与常数1比较，积分发散。这里没有使用尚待建立的指数幂级数展开；
也没有把“被积函数趋零”当作一般反常积分收敛的必要条件。
\end{solution}

\begin{solution}{improper-log-power-zero}
令 \(t=-\log x\)，则
\[
 \int_0^{1/2}\frac{dx}{x^p\abs{\log x}^{q}}
 =\int_{\log2}^{\infty}e^{(p-1)t}t^{-q}\,dt.
\]
因此 \(p<1\) 时收敛，\(p>1\) 时发散；\(p=1\) 时恰在 \(q>1\)
时收敛。
\end{solution}

\begin{solution}{improper-sine-over-power}
当 \(\alpha>0\) 时，\(x^{-\alpha}\) 单调趋零，Dirichlet 判别给出收敛。
绝对收敛在 \(\alpha>1\) 时由 \(p\)-型比较得到。若
\(\alpha\le1\)，则在 \(x\ge1\) 有
\(\abs{\sin x}/x^\alpha\ge\abs{\sin x}/x\)。正文已用有限和与
\(\int dt/t\) 的比较证明后者积分发散，故绝对积分也发散。

当 \(\alpha=0\) 时截断积分 \(\cos1-\cos R\) 无极限；当
\(\alpha<0\) 时令 \(\beta=-\alpha>0\)，并取
\[
 u_k=2k\pi+\frac{\pi}{6},\qquad
 v_k=2k\pi+\frac{5\pi}{6}.
\]
在 \([u_k,v_k]\) 上 \(\sin x\ge1/2\)，故
\[
 \int_{u_k}^{v_k}x^\beta\sin x\,dx
 \ge\frac{v_k-u_k}{2}\,u_k^\beta\longrightarrow+\infty.
\]
这给出两端都趋于无穷而尾积分不趋零的区间列，Cauchy 准则失败。
所以 \(\alpha>1\) 时绝对收敛，\(0<\alpha\le1\) 时条件收敛，
\(\alpha\le0\) 时发散。
\end{solution}

\begin{solution}{improper-dirichlet-quantitative-tail}
对 \(v>u\) 的分部积分估计给出
\[
 \abs{\int_u^vfg}\le2M\abs{g(u)}.
\]
令 \(v\to\infty\)，左端趋于尾积分，右端不变，故得到所需估计。
\end{solution}

\begin{solution}{improper-abel-from-dirichlet}
单调有界的 \(g\) 有极限 \(L\)。由 \(\int f\) 收敛，函数
\(F(x)=\int_a^x f\) 在无穷远有有限极限，因此在某个尾部有界；在
余下有限闭区间上，它的连续性保证有界。第一项 \(Lf\) 的反常积分
收敛；第二项 \(f(g-L)\) 中，\(f\) 连续、\(g-L\) 连续可导且单调趋零，满足
Dirichlet 判别的全部条件。两项之和收敛。
\end{solution}

\begin{solution}{improper-absolute-not-necessary}
收敛性由 Dirichlet 判别。绝对值积分发散的证明如下：对充分大的 \(k\)，
\(I_k=[k\pi+\pi/6,k\pi+5\pi/6]\) 上
\(\abs{\sin x}\ge1/2\)，且 \(x\le(k+1)\pi\)，所以
\[
 \int_{I_k}\frac{\abs{\sin x}}x\,dx
 \ge \frac{1}{2(k+1)\pi}\cdot\frac{2\pi}{3}
 =\frac1{3(k+1)}.
\]
对 \(k=1,\ldots,N\) 求和，所得下界至少为
\(\frac13\int_2^{N+2}dt/t=\frac13\log((N+2)/2)\)，趋于正无穷。
因此非负的截断积分无界，不绝对收敛。
\end{solution}

\begin{solution}{improper-substitution-inversion}
先对 \(0<\varepsilon<1\) 作有限区间代换 \(x=1/t\)，有
\[
 \int_\varepsilon^1 f(x)\,dx
 =\int_1^{1/\varepsilon}\frac{f(1/t)}{t^2}\,dt.
\]
\(\varepsilon\downarrow0\) 与 \(1/\varepsilon\to\infty\) 一一对应，
所以两边的反常积分同敛散，收敛时相等。特别地，取 \(f(x)=x^{-p}\)，得
\[
 \int_0^1 x^{-p}\,dx
 =\int_1^\infty t^{p-2}\,dt
 =\int_1^\infty\frac{dt}{t^{2-p}}.
\]
右端在 \(2-p>1\)，即 \(p<1\) 时收敛，与零点附近的判别一致。
\end{solution}

\begin{solution}{improper-gaussian-convergence}
在 \([0,1]\) 上函数连续。对 \(x\ge1\)，有 \(x^2\ge x\)，故
\[
 0\le e^{-x^2}\le e^{-x}.
\]
而 \(\int_1^\infty e^{-x}dx\) 收敛，比较判别给出结论。
\end{solution}

\begin{solution}{improper-integration-by-parts-boundary}
取 \(u(x)\equiv1\)、\(v(x)=\sin x\)。则
\[
 \int_0^\infty u'(x)v(x)\,dx=0
\]
收敛，但 \(u(R)v(R)=\sin R\) 无极限。相应的
\(\int_0^\infty uv'=\int_0^\infty\cos x\,dx\) 也发散，说明缺少边界项
极限时不能由形式分部积分推出收敛。
\end{solution}

\begin{solution}{improper-conditional-substitution-warning}
若 \(\phi\) 是题设类型的单调 \(C^1\) 满射，则每个有限截断满足
\[
 \int_{\alpha}^{s}
 \frac{\sin(\phi(t))}{\phi(t)}\phi'(t)\,dt
 =\int_1^{\phi(s)}\frac{\sin x}{x}\,dx.
\]
令 \(s\) 趋于端点，右端按原顺序趋于原积分值。任意重排则把远处不同
区间的积分块改变顺序，并不由单调变量代换产生；条件收敛时这种重排可能
改变极限。
\end{solution}

\begin{solution}{improper-fresnel-dirichlet}
在 \([1,\infty)\) 令 \(t=x^2\)，得到
\[
 \int_1^\infty\sin(x^2)\,dx
 =\frac12\int_1^\infty\frac{\sin t}{\sqrt t}\,dt,
\]
由 Dirichlet 判别收敛，\([0,1]\) 上连续因而无问题。

若取绝对值，则
\[
 \frac12\int_1^\infty\frac{\abs{\sin t}}{\sqrt t}\,dt
\]
大于等于 \(\frac12\int_1^\infty\abs{\sin t}/t\,dt\)，后者已经用
有限截断比较证明发散。因此原积分不绝对收敛。
\end{solution}
